%Data tentang institusi

%\input{Data guru}

%\renewcommand{\fachbereich}{Departemen}

%\renewcommand{\dozent}{Prof. Dr. }

%Tanggal ujian

\renewcommand{\klausurgebiet}{ }

\renewcommand{\klausurtyp}{ }

\renewcommand{\klausurdatum}{. 20}

\klausurvorspann {\fachbereich} {\klausurdatum} {\dozent} {\klausurgebiet} {\klausurtyp}

%Data untuk tabel poin berikut

\renewcommand{\aeins}{  3  }

\renewcommand{\azwei}{  3   }

\renewcommand{\adrei}{  6  }

\renewcommand{\avier}{  5  }

\renewcommand{\afuenf}{  5  }

\renewcommand{\asechs}{  2  }

\renewcommand{\asieben}{  5  }

\renewcommand{\aacht}{  5  }

\renewcommand{\aneun}{  7  }

\renewcommand{\azehn}{  3  }

\renewcommand{\aelf}{  2  }

\renewcommand{\azwoelf}{  5  }

\renewcommand{\adreizehn}{  9   }

\renewcommand{\avierzehn}{  5  }

\renewcommand{\afuenfzehn}{  65  }

\renewcommand{\asechzehn}{    }

\renewcommand{\asiebzehn}{    }

\renewcommand{\aachtzehn}{    }

\renewcommand{\aneunzehn}{    }

\renewcommand{\azwanzig}{    }

\renewcommand{\aeinundzwanzig}{    }

\renewcommand{\azweiundzwanzig}{    }

\renewcommand{\adreiundzwanzig}{    }

\renewcommand{\avierundzwanzig}{    }

\renewcommand{\afuenfundzwanzig}{    }

\renewcommand{\asechsundzwanzig}{    }

\punktetabellevierzehn

\klausurnote

\newpage

\setcounter{section}{0}



\inputaufgabeklausurloesung
{3}

{

Definisikan istilah-istilah berikut
\zusatzklammer {yang dicetak miring} {} {}.
\aufzaehlungsechs{\stichwort {Lingkaran kelengkungan} {}
dari suatu
\definitionsverweis {kurva}{}{}
yang dua kali
\definitionsverweis {terdiferensialkan secara kontinu}{}{}
dan
\definitionsverweis {berparametrisasi panjang busur}{}{}.
Misalkan
\maabbdisp {\gamma} {I} {\R^2
} {}
dan ambil titik

\mavergleichskette
{\vergleichskette
{ t_0
}
{ \in }{ I
}
{  }{
}
{    }{
}
{   }{
}
}
{}{}{.}

}{Suatu \stichwort {submanifold tertutup} {}
\mathl{M \subseteq N}{} dari manifold diferensiabel $N$.

}{\stichwort {Ruang singgung} {} pada titik
\mathl{P\in M}{} dari manifold diferensiabel $M$.

}{Istilah
\stichwort {berorientasi} {}
untuk suatu
\definitionsverweis {manifold diferensiabel}{}{}
$M$.

}{\stichwort {Batas} {}
suatu
\definitionsverweis {manifold berbatas}{}{.}

}{\stichwort {Simbol Christoffel} {}

\mathl{\Gamma_{ij}^k}{} bagi koneksi Levi-Civita dari struktur Riemann $\left(  g_{ij}  \right)$ pada himpunan terbuka

\mavergleichskette
{\vergleichskette
{ U
}
{ \subseteq }{ \R^n
}
{  }{
}
{    }{
}
{   }{
}
}
{}{}{.}
}

}
{

\aufzaehlungsechs{Lingkaran dengan jari-jari

\mavergleichskettedisp
{\vergleichskette
{ r
}
{ =} {  { \frac{  1 }{ \betrag {  \gamma_1'(t_0) \gamma_2^{\prime \prime} (t_0) - \gamma_2'(t_0) \gamma_1^{\prime \prime} (t_0)  }  } }
}
{ } {
}
{ } {
}
{ } {
}
}
{}{}{}
dan pusat

\mavergleichskettedisp
{\vergleichskette
{ M
}
{ =} { \gamma(t_0) +  { \frac{  1 }{  \gamma_1'(t_0) \gamma_2^{\prime \prime} (t_0) - \gamma_2'(t_0) \gamma_1^{\prime \prime} (t_0)  } }  \begin{pmatrix} - \gamma_2'(t_0) \\ \gamma_1'(t_0)   \end{pmatrix}
}
{ } {
}
{ } {
}
{ } {
}
}
{}{}{}
dinamakan lingkaran kelengkungan dari $\gamma$ di $t_0$.
}{Suatu himpunan bagian tertutup
\mathl{M \subseteq N}{} dinamakan submanifold tertutup jika untuk setiap titik
\mathl{P \in M}{} terdapat suatu
\definitionsverweis {peta}{}{}
sedemikian sehingga
\mathl{P \in W \subseteq N}{},
\definitionsverweis {terbuka}{}{,}
\maabb {\theta} { W } { W'
} {,}

\mathl{W' \subseteq \R^n}{} terbuka, dan

\mavergleichskettedisp
{\vergleichskette
{ M \cap W
}
{ =} {\theta^{-1} (( \R^m \times\{0\}) \cap W' )
}
{ } {
}
{ } {
}
{ } {
}
}
{}{}{.}
}{Ruang singgung
\mathl{T_PM}{} terdiri atas semua kelas ekuivalensi lintasan terdiferensialkan yang melalui titik tersebut dan ekuivalen secara tangensial.
}{Suatu
\definitionsverweis {manifold diferensiabel}{}{}
$M$ dengan sebuah
\definitionsverweis {atlas}{}{}

\mathl{(U_i, V_i, \alpha_i)}{} disebut berorientasi jika setiap petanya
\definitionsverweis {berorientasi}{}{}
dan semua pergantian petanya
\definitionsverweis {mempertahankan orientasi}{}{}
.
}{Batas $M$ didefinisikan oleh

\mavergleichskettedisp
{\vergleichskette
{ \partial M
}
{ =} {  \left\{ x \in M  \mid   \alpha_i(x) \in  \partial H \cap V_i  \text{ untuk } \text{suatu } i \in I         \right\}
}
{ } {
}
{ } {
}
{ } {
}
}
{}{}{}
di mana $\alpha_i$ adalah peta-peta koordinat.
}{Misalkan
\mathl{\left(  g^{k l}  \right)}{} adalah matriks invers dari
\mathl{\left(  g_{ij}  \right)}{.}
\zusatzklammer {Fungsi-fungsi bernilai riil pada $U$} {} {}

\mavergleichskettedisp
{\vergleichskette
{\Gamma_{ij}^k
}
{ =} {  { \frac{  1 }{ 2 } } \sum_{l = 1}^n g^{kl}(\partial_ig_{jl}+\partial_jg_{il}-\partial_lg_{ij})
}
{ } {
}
{ } {
}
{ } {
}
}
{}{}{}
dinamakan simbol-simbol Christoffel untuk koneksi Levi-Civita.
}

 }



\inputaufgabeklausurloesung
{3}

{

Nyatakan teorema-teorema berikut.
\aufzaehlungdrei{Teorema tentang citra
\definitionsverweis {pemetaan Gauss}{}{}
dari suatu hipermuka.}{Rumus luas permukaan dari permukaan putar yang dibentuk oleh
\definitionsverweis {kurva diferensiabel}{}{}
\maabbeledisp {\gamma} {]a,b[} {\R^2
} { t } {(x(t),y(t))
} {,}
dengan

\mavergleichskette
{\vergleichskette
{  y(t)
}
{ > }{  0
}
{  }{
}
{    }{
}
{   }{
}
}
{}{}{.}}{\stichwort {Aturan hasil kali} {}
untuk turunan eksterior.}

}
{

\aufzaehlungdrei{\faktsituation {Misalkan

\mavergleichskette
{\vergleichskette
{ W
}
{ \subseteq }{\R^n
}
{ }{
}
{ }{
}
{ }{
}
}
{}{}{}
adalah suatu
\definitionsverweis {himpunan bagian terbuka}{}{}
dan
\maabbdisp {h} { W } {\R
} {}
adalah suatu
\definitionsverweis {fungsi terdiferensialkan secara kontinu}{}{.}}
\faktvoraussetzung {
\definitionsverweis {serat}{}{}

\mavergleichskette
{\vergleichskette
{ Y
}
{ = }{ h^{-1}(c)
}
{ }{
}
{ }{
}
{ }{
}
}
{}{}{}
dari $h$ di

\mavergleichskette
{\vergleichskette
{ c
}
{ \in }{\R
}
{  }{
}
{    }{
}
{   }{
}
}
{}{}{}
bersifat
\definitionsverweis {kompakt}{}{}
dan
\definitionsverweis {reguler}{}{.}}
\faktfolgerung {Maka
\definitionsverweis {pemetaan Gauss}{}{}
\maabbdisp {} { Y } { S^{n-1}
} {}
bersifat surjektif.}
\faktzusatz {}
\faktzusatz {}}{\faktsituation {Misalkan
\maabbeledisp {\gamma} {]a,b[} {\R^2
} { t } {(x(t),y(t))
} {,}
adalah suatu
\definitionsverweis {kurva terdiferensialkan}{}{}
dengan

\mavergleichskette
{\vergleichskette
{ y(t)
}
{ > }{ 0
}
{ }{
}
{ }{
}
{ }{
}
}
{}{}{,}}
\faktvoraussetzung {yang menginduksi suatu
\definitionsverweis {difeomorfisme}{}{}
ke

\mavergleichskette
{\vergleichskette
{M
}
{ = }{ \gamma(I)
}
{ }{
}
{ }{
}
{ }{
}
}
{}{}{}
, dengan

\mavergleichskette
{\vergleichskette
{M
}
{ \subseteq }{ G
}
{ }{
}
{ }{
}
{ }{
}
}
{}{}{}
sebagai
\definitionsverweis {submanifold tertutup}{}{}
satu dimensi dalam suatu himpunan terbuka

\mavergleichskette
{\vergleichskette
{G
}
{ \subseteq }{\R \times \R_+
}
{ }{
}
{ }{
}
{ }{
}
}
{}{}{}
.}
\faktfolgerung {Maka
\definitionsverweis {permukaan putar}{}{}
yang bersesuaian merupakan submanifold tertutup dari
\mathl{\R^3}{} tanpa sumbu $x$, dan luas permukaannya adalah

\mathdisp {2 \pi \int_{  a  }^{  b  } \sqrt{ (x'(t))^2+ (y'(t))^2 } \cdot y(t) \, d t} { . }
}
\faktzusatz {}
\faktzusatz {}}{\faktsituation {Misalkan

\mavergleichskette
{\vergleichskette
{ U
}
{ \subseteq }{ \R^n
}
{ }{
}
{ }{
}
{ }{
}
}
{}{}{}
\definitionsverweis {terbuka}{}{}
dan misalkan

\mavergleichskette
{\vergleichskette
{ \omega
}
{ \in }{
{ \mathcal E }^{  k  }_{ 1 } (  U   )
}
{  }{
}
{    }{
}
{   }{
}
}
{}{}{}
serta

\mavergleichskette
{\vergleichskette
{ \tau
}
{ \in }{
{ \mathcal E }^{  \ell }_{ 1 } (  U   )
}
{  }{
}
{    }{
}
{   }{
}
}
{}{}{}
adalah bentuk-bentuk diferensial terdiferensialkan pada $U$.}
\faktfolgerung {Maka berlaku aturan hasil kali

\mavergleichskettedisp
{\vergleichskette
{ d( \omega \wedge \tau)
}
{ =} { ( d \omega) \wedge \tau  + (-1)^k \omega  \wedge  (d \tau)
}
{ } {
}
{ } {
}
{ } {
}
}
{}{}{.}}
\faktzusatz {}
\faktzusatz {}}

 }



\inputaufgabeklausurloesung
{6 (1+2+3)}

{

Misalkan $Y$ adalah hipermuka yang diberikan oleh persamaan

\mavergleichskettedisp
{\vergleichskette
{ 2x^2 +3y^2+5z^2
}
{ =} { 10
}
{ } {
}
{ } {
}
{ } {
}
}
{}{}{}
di $\R^3$. Misalkan

\mavergleichskette
{\vergleichskette
{ P
}
{ = }{ \begin{pmatrix}  -1 \\ -1\\  1   \end{pmatrix}
}
{ \in }{ Y
}
{    }{
}
{   }{
}
}
{}{}{}
dan

\mavergleichskette
{\vergleichskette
{ w
}
{ = }{ \begin{pmatrix}  -3 \\ 2 \\  4   \end{pmatrix}
}
{ \in }{ \R^3
}
{    }{
}
{   }{
}
}
{}{}{.}
\aufzaehlungdrei{Tentukan
\definitionsverweis {ruang singgung}{}{}

\mathl{T_PY}{.}
}{Tentukan
\definitionsverweis {dekomposisi ortogonal}{}{}
$w$ menjadi komponen tangensial dan komponen ortogonal pada titik $P$.
}{Realisasikan komponen tangensial $w$ di $T_PY$ melalui suatu kurva diferensiabel yang seluruhnya terletak pada $Y$.
}

}
{

\aufzaehlungdrei{Total diferensialnya adalah

\mathdisp {\left(  4x  , \,   6y  , \,   10z                 \right)} { , }

sehingga pada $P$ diperoleh

\mavergleichskettedisp
{\vergleichskette
{ T_PY
}
{ =} { \left\{  \begin{pmatrix}  a  \\ b \\ c   \end{pmatrix}   \mid   -4a-6b+10c = 0         \right\}
}
{ =} { \left\{  \begin{pmatrix}  a  \\ b \\ c   \end{pmatrix}   \mid      \left\langle  \begin{pmatrix}  a  \\ b \\ c   \end{pmatrix}   , \begin{pmatrix}  -4 \\ -6\\  10   \end{pmatrix}     \right\rangle = 0         \right\}
}
{ } {
}
{ } {
}
}
{}{}{.}
}{Kita gunakan ansatz

\mavergleichskettedisp
{\vergleichskette
{  \begin{pmatrix}  -3 \\ 2\\  4   \end{pmatrix}
}
{ =} { s  \begin{pmatrix}  -4 \\ -6\\  10   \end{pmatrix} + v
}
{ } {
}
{ } {
}
{ } {
}
}
{}{}{}
dengan skalar yang dicari

\mavergleichskette
{\vergleichskette
{ s
}
{ \in }{ \R
}
{ }{
}
{ }{
}
{ }{
}
}
{}{}{}
dan vektor tangensial yang dicari $v$. Kita peroleh

\mavergleichskettedisp
{\vergleichskette
{ \left\langle \begin{pmatrix} -3 \\ 2\\ 4 \end{pmatrix} , \begin{pmatrix} -4 \\ -6\\ 10 \end{pmatrix} \right\rangle
}
{ =} { 12-12+40
}
{ =} { 40
}
{ } {
}
{ } {
}
}
{}{}{.}
Karena

\mavergleichskettedisp
{\vergleichskette
{ \left\langle  \begin{pmatrix}  -4 \\ -6\\  10   \end{pmatrix}  ,  \begin{pmatrix}  -4 \\ -6\\  10   \end{pmatrix}   \right\rangle
}
{ =} { 16 +36 + 100
}
{ =} { 152
}
{ } {
}
{ } {
}
}
{}{}{}
maka

\mavergleichskettedisp
{\vergleichskette
{ s
}
{ =} { { \frac{   40 }{  152 } }
}
{ =} { { \frac{   5 }{  19 } }
}
{ } {
}
{ } {
}
}
{}{}{}
dan dengan demikian

\mavergleichskettedisp
{\vergleichskette
{ v
}
{ =} { \begin{pmatrix}  -3 \\ 2\\  4   \end{pmatrix}   - { \frac{   5 }{  19 } }  \begin{pmatrix}  -4 \\ -6\\  10   \end{pmatrix}
}
{ =} { { \frac{   1  }{  19 } }  \begin{pmatrix}  -57+20 \\ 38 +30\\  76 -50   \end{pmatrix}
}
{ =} { { \frac{   1  }{  19 } }  \begin{pmatrix}  -37 \\ 68\\  26   \end{pmatrix}
}
{ } {
}
}
{}{}{.}
Jadi, dekomposisi ortogonalnya adalah

\mavergleichskettedisp
{\vergleichskette
{  w
}
{ =} { \begin{pmatrix}  -3 \\ 2\\  4   \end{pmatrix}
}
{ =} { { \frac{   5 }{  19 } }  \begin{pmatrix}  -4 \\ -6\\  10   \end{pmatrix} + { \frac{   1  }{  19 } }  \begin{pmatrix}  -37 \\ 68\\  26   \end{pmatrix}
}
{ } {
}
{ } {
}
}
{}{}{.}
}{Kita tulis persamaan permukaan

\mavergleichskette
{\vergleichskette
{ 2x^2 +3y^2+5z^2
}
{ = }{ 10
}
{ }{
}
{ }{
}
{ }{
}
}
{}{}{}
sebagai

\mavergleichskettedisp
{\vergleichskette
{ z
}
{ =} { \pm \sqrt{ 10- 2x^2 - 3y^2 }
}
{ } {
}
{ } {
}
{ } {
}
}
{}{}{,}
di mana akar positif harus dipilih pada titik $P$. Ini merupakan parameterisasi permukaan pada suatu lingkungan terbuka dari $P$; ruang singgung di
\mathl{\begin{pmatrix} -1 \\ -1 \end{pmatrix}}{} dipetakan ke ruang singgung
\mathl{T_PY}{}.
}

 }



\inputaufgabeklausurloesung
{5}

{

Misalkan
\maabbeledisp {\gamma} {\R} { \R^2
} { t } { \begin{pmatrix}  \cos \left(  \omega t \right)  \\ \sin  \left(  \omega t \right)    \end{pmatrix}
} {,}
dengan

\mavergleichskette
{\vergleichskette
{ \omega
}
{ > }{ 0
}
{  }{
}
{    }{
}
{   }{
}
}
{}{}{.}
Tunjukkan bahwa tidak ada gerak melingkar lain dengan norma kecepatan konstan yang memiliki nilai dan turunan hingga orde kedua yang sama dengan $\gamma$ di titik parameter

\mavergleichskette
{\vergleichskette
{ t
}
{ = }{ 0
}
{  }{
}
{    }{
}
{   }{
}
}
{}{}{}
tersebut.

}
{

Kita peroleh

\mavergleichskettedisp
{\vergleichskette
{ \gamma(0)
}
{ =} { \begin{pmatrix} 1 \\ 0 \end{pmatrix}
}
{ } {
}
{ } {
}
{ } {
}
}
{}{}{,}

\mavergleichskettedisp
{\vergleichskette
{ \gamma'(0)
}
{ =} { \begin{pmatrix} 0 \\ \omega \end{pmatrix}
}
{ } {
}
{ } {
}
{ } {
}
}
{}{}{}
dan

\mavergleichskettedisp
{\vergleichskette
{ \gamma^{\prime \prime} (0)
}
{ =} { \begin{pmatrix} - \omega^2 \\ 0 \end{pmatrix}
}
{ } {
}
{ } {
}
{ } {
}
}
{}{}{.}
Suatu gerak melingkar yang sama dengan gerak yang diberikan pada

\mavergleichskette
{\vergleichskette
{ t
}
{ = }{ 0
}
{ }{
}
{ }{
}
{ }{
}
}
{}{}{}
hingga turunan kedua khususnya harus memiliki garis singgung yang sama di titik tersebut. Oleh karena itu, pusat lingkarannya harus terletak pada sumbu $x$. Selain itu, pusat itu harus terletak di sebelah kiri
\mathl{\begin{pmatrix} 1 \\ 0 \end{pmatrix}}{} karena percepatannya mengarah ke kiri. Karena itu, gerak melingkar beraturan yang dicari dengan pusat
\mathl{\begin{pmatrix} x \\ 0 \end{pmatrix}}{}
\zusatzklammer {
\mavergleichskettek
{\vergleichskettek
{ x
}
{ < }{ 1
}
{ }{
}
{ }{
}
{ }{
}
}
{}{}{}} {} {}
dan jari-jari
\mathl{1-x}{} dapat dituliskan sebagai

\mavergleichskettedisp
{\vergleichskette
{  \delta (t)
}
{ =} {  \begin{pmatrix}  x  \\ 0   \end{pmatrix} + (1-x) \begin{pmatrix}  \cos \left(  u t \right)  \\ \sin  \left(  u t \right)     \end{pmatrix}
}
{ } {
}
{ } {
}
{ } {
}
}
{}{}{}
. Kita mempunyai

\mavergleichskettedisp
{\vergleichskette
{ \delta(0)
}
{ =} { \gamma(0)
}
{ } {
}
{ } {
}
{ } {
}
}
{}{}{.}
Selanjutnya,

\mavergleichskettedisp
{\vergleichskette
{ \delta'( 0)
}
{ =} { (1-x) \begin{pmatrix} 0 \\ u \end{pmatrix}
}
{ } {
}
{ } {
}
{ } {
}
}
{}{}{}
dan

\mavergleichskettedisp
{\vergleichskette
{ \delta^{\prime \prime}( 0)
}
{ =} { (1-x) \begin{pmatrix} -u^2 \\ 0 \end{pmatrix}
}
{ } {
}
{ } {
}
{ } {
}
}
{}{}{.}
Agar syarat-syarat yang diminta terpenuhi, harus berlaku

\mavergleichskettedisp
{\vergleichskette
{ \omega
}
{ =} { (1-x) u
}
{ } {
}
{ } {
}
{ } {
}
}
{}{}{}
dan

\mavergleichskettedisp
{\vergleichskette
{ \omega^2
}
{ =} { (1-x) u^2
}
{ } {
}
{ } {
}
{ } {
}
}
{}{}{}
. Dari sini diperoleh

\mavergleichskettedisp
{\vergleichskette
{ \omega^2
}
{ =} { \omega u
}
{ } {
}
{ } {
}
{ } {
}
}
{}{}{}
dan dengan demikian

\mavergleichskettedisp
{\vergleichskette
{ u
}
{ =} { \omega
}
{ } {
}
{ } {
}
{ } {
}
}
{}{}{}
serta

\mavergleichskettedisp
{\vergleichskette
{ x
}
{ =} { 0
}
{ } {
}
{ } {
}
{ } {
}
}
{}{}{.}

 }



\inputaufgabeklausurloesung
{5}

{

Buktikan teorema tentang kelengkungan suatu kurva bidang yang diberikan secara implisit.

}
{

Misalkan
\maabb {\gamma} { I } { \R^2
} {}
adalah parameterisasi panjang busur kurva $Y$ pada suatu lingkungan dari $P$

\mavergleichskette
{\vergleichskette
{ \gamma (0)
}
{ = }{ P
}
{ }{
}
{ }{
}
{ }{
}
}
{}{}{,}
yang sesuai dengan orientasi yang diberikan. Diferensial total
\maabbdisp {\left(D N \right)_{ P }} { \R^2} {\R^2
} {}
bersifat linier, sehingga cukup membuktikan pernyataan bagi vektor $\gamma'(0)$. Berdasarkan
aturan rantai,

\mavergleichskettedisp
{\vergleichskette
{ \left(  D N   \right)_{ P } \left(   \gamma'(0)  \right)
}
{ =} { (N \circ \gamma )'(0)
}
{ } {
}
{ } {
}
{ } {
}
}
{}{}{.}
Vektor ini merupakan kelipatan dari $\gamma'(0)$, sehingga sama dengan

\mathdisp {\left\langle  (N \circ  \gamma)'(0) ,  \gamma' (0)   \right\rangle  \gamma'(0)} { . }

Karena syarat ortogonalitas memberikan

\mavergleichskettedisp
{\vergleichskette
{ \left\langle (N \circ \gamma)(t) , \gamma' (t) \right\rangle
}
{ =} { 0
}
{ } {
}
{ } {
}
{ } {
}
}
{}{}{}
untuk semua

\mavergleichskette
{\vergleichskette
{ t
}
{ \in }{ I
}
{ }{
}
{ }{
}
{ }{
}
}
{}{}{}
dan karenanya

\mavergleichskettedisp
{\vergleichskette
{ \left\langle (N \circ \gamma)'(t) , \gamma' (t) \right\rangle + \left\langle (N \circ \gamma)(t) , \gamma^{\prime \prime} (t) \right\rangle
}
{ =} { 0
}
{ } {
}
{ } {
}
{ } {
}
}
{}{}{.}
Dengan demikian,

\mavergleichskettedisp
{\vergleichskette
{ \left(  D N   \right)_{ P } \left(   \gamma'(0)  \right)
}
{ =} { - \left\langle  (N \circ  \gamma)(0) ,  \gamma^{\prime \prime} (0)   \right\rangle  \gamma'(0)
}
{ =} { - \kappa( \gamma(0)) \gamma'(0)
}
{ } {
}
{ } {
}
}
{}{}{}
berdasarkan
Lemma 3.8 (Differentialgeometrie (Osnabrück 2023)).

 }



\inputaufgabeklausurloesung
{2}

{

Berikan contoh manifold diferensiabel
\mathkor {} {M} {dan} {N} {}
dengan

\mavergleichskette
{\vergleichskette
{ \operatorname{ dim}_{  } ^{  }  \left(   M   \right), \, \operatorname{ dim}_{  } ^{  }  \left(   N   \right)
}
{ \geq }{ 1
}
{  }{
}
{    }{
}
{   }{
}
}
{}{}{}
serta pemetaan diferensiabel
\maabb {\varphi} { M } { N
} {}
dan
\maabb {\psi} { N } { M
} {}
sedemikian sehingga

\mavergleichskette
{\vergleichskette
{ \psi \circ \varphi
}
{ = }{
\operatorname{Id}_{ M }
}
{  }{
}
{    }{
}
{   }{
}
}
{}{}{}
dan

\mavergleichskette
{\vergleichskette
{\varphi \circ  \psi
}
{ \neq }{
\operatorname{Id}_{ N }
}
{  }{
}
{    }{
}
{   }{
}
}
{}{}{}
berlaku.

}
{

Kita pilih

\mavergleichskette
{\vergleichskette
{M
}
{ = }{\R
}
{ }{
}
{ }{
}
{ }{
}
}
{}{}{}
dan

\mavergleichskette
{\vergleichskette
{N
}
{ = }{\R^2
}
{ }{
}
{ }{
}
{ }{
}
}
{}{}{}
dan
\maabbeledisp {\varphi} {\R} {\R^2
} { x } {(x,0)
} {,}
dan
\maabbeledisp {\psi} {\R^2} {\R
} {(x,y)} { x
} {.}
Maka

\mavergleichskettedisp
{\vergleichskette
{ \left( \psi \circ \varphi \right) (x)
}
{ =} { \psi((x,0))
}
{ =} { x
}
{ } {
}
{ } {
}
}
{}{}{}
dan

\mavergleichskettedisp
{\vergleichskette
{ \left(  \varphi  \circ \psi  \right) ((x,y))
}
{ =} { \varphi (x)
}
{ =} {(x,0)
}
{ } {
}
{ } {
}
}
{}{}{.}

 }



\inputaufgabeklausurloesung
{5}

{

Buktikan bahwa ekuivalensi tangensial lintasan-lintasan pada suatu manifold di titik $P$ dapat diperiksa dengan menggunakan sebarang peta koordinat.

}
{

Untuk suatu
\definitionsverweis {kurva terdiferensialkan}{}{}
\maabbdisp {\gamma} { I } { M
} {}
dengan

\mavergleichskette
{\vergleichskette
{ 0
}
{ \in }{ I
}
{ }{
}
{ }{
}
{ }{
}
}
{}{}{}
dan

\mavergleichskette
{\vergleichskette
{ \gamma(0)
}
{ = }{ P
}
{  }{
}
{    }{
}
{   }{
}
}
{}{}{}
dan suatu
\definitionsverweis {peta}{}{}
\maabbdisp {\alpha} { U } { V
} {}
\zusatzklammer {dengan

\mavergleichskettek
{\vergleichskettek
{P
}
{ \in }{ U
}
{ }{
}
{ }{
}
{ }{
}
}
{}{}{}
dan

\mavergleichskette
{\vergleichskette
{V
}
{ \subseteq }{\R^n
}
{ }{
}
{ }{
}
{ }{
}
}
{}{}{}} {} {}
ekspresi

\mathdisp {\left(  \alpha \circ \left(  \gamma \vert_{\gamma^{-1} (U)}  \right)  \right)'(0)} {  }

tidak berubah apabila kita beralih ke interval terbuka yang lebih kecil

\mavergleichskette
{\vergleichskette
{ 0
}
{ \in }{ I'
}
{ \subseteq }{ I
}
{ }{
}
{ }{
}
}
{}{}{}
dan himpunan terbuka yang lebih kecil

\mavergleichskette
{\vergleichskette
{P
}
{ \in }{ U'
}
{ \subseteq }{ U
}
{ }{
}
{ }{
}
}
{}{}{}
\zusatzklammer {dengan peta yang diinduksi} {} {}
. Karena itu, kita dapat mengasumsikan bahwa
\mathkor {} {\gamma_1} {dan} {\gamma_2} {}
didefinisikan pada interval yang sama, citranya terletak di $U$, dan pada $U$ terdapat dua peta
\maabbdisp {\alpha_1} { U } { V_1
} {}
dan
\maabbdisp {\alpha_2} { U } { V_2
} {}
. Selanjutnya, dari

\mavergleichskettedisp
{\vergleichskette
{ ( \alpha_1 \circ \gamma_1 )'(0)
}
{ =} { ( \alpha_1 \circ \gamma_2 )'(0)
}
{ } {
}
{ } {
}
{ } {
}
}
{}{}{}
berdasarkan
aturan rantai
dan keterdiferensialan
\definitionsverweis {pemetaan transisi}{}{}

\mathl{\alpha_2 \circ \alpha_1^{-1}}{}, langsung diperoleh

\mavergleichskettealign
{\vergleichskettealign
{ \left(  \alpha_2 \circ \gamma_1  \right)'(0)
}
{ =} { \left(  \left(  \alpha_2 \circ \alpha_1^{-1}  \right) \circ \left(  \alpha_1 \circ \gamma_1  \right)  \right)'(0)
}
{ =} { \left(  D   \left(   \alpha_2 \circ \alpha_1^{-1}   \right)      \right)_{\alpha_1(P)} \left(   \left(  \alpha_1 \circ \gamma_1  \right)'(0)   \right)
}
{ =} { \left(  D   \left(   \alpha_2 \circ \alpha_1^{-1}   \right)      \right)_{\alpha_1(P)} \left(   \left(  \alpha_1 \circ \gamma_2  \right)'(0)   \right)
}
{ =} { \left(  \alpha_2 \circ \gamma_2  \right)'(0)
}
}
{}
{}{.}

 }



\inputaufgabeklausurloesung
{5}

{

Tunjukkan bahwa pemetaan singgung $T(\varphi)$ dari
\maabbeledisp {\varphi} {\R^1} {S^1
} { t } {( \cos  t, \sin  t  )
} {,}
bersifat surjektif.

}
{

Pemetaan $\varphi$ bersifat surjektif, sehingga cukup ditunjukkan bahwa untuk setiap
\mathl{t \in \R}{} pemetaan singgung linier
\maabbdisp {} {T_t \R \cong \R} { T_{\varphi(t)} S^1
} {}
bersifat surjektif. Karena kedua ruang itu berdimensi satu, cukup ditunjukkan bahwa suatu vektor yang berbeda dari $0$ tidak dipetakan ke $0$. Vektor singgung di $t$ direalisasikan oleh lintasan terdiferensialkan
\maabbeledisp {\gamma} {\R} {\R
} { s } {t+s
} {.}
Lintasan komposisi
\maabbeledisp {\varphi \circ \gamma} {\R} {S^1
} { s } { ( \cos (t+s) , \sin (t+s) )
} {,}
mewujudkan vektor singgung citra; pada bidang ambien,
\zusatzklammer {
\mathl{T_{\varphi(t)}S^1 \subseteq \R^2}{}} {} {}

\mavergleichskettedisp
{\vergleichskette
{( \varphi \circ \gamma)' (0)
}
{ =} {(- \sin t , \cos t )
}
{ } {
}
{ } {
}
{ } {
}
}
{}{}{,}
dan vektor ini tidak nol.

 }



\inputaufgabeklausurloesung
{7 (3+4)}

{

Tinjau pemetaan diferensiabel
\maabbeledisp {\gamma} { [1,c] } {\R^2
} { t } {(t,t^{3})
} {,}
\zusatzklammer {dengan \mathlk{c \geq 1}{}} {} {}
dan
\maabbeledisp {\varphi} {\R_+ \times \R_+ } {\R^3
} {(u,v)} {(u^3,u^2+v^2,u^{-1}v^{-1} )
} {,}
serta bentuk diferensial

\mavergleichskettedisp
{\vergleichskette
{ \omega
}
{ =} { (x-y)dx-z^2dy+dz
}
{ } {
}
{ } {
}
{ } {
}
}
{}{}{}
pada $\R^3$.

a) Hitung bentuk diferensial tarik balik
\mathl{\varphi^* \omega}{} pada $\R_+ \times \R_+$.

b) Hitung integral lintasan dari bentuk diferensial
\mathl{\varphi^* \omega}{} sepanjang lintasan $\gamma$ sebagai fungsi dari $c$.

}
{

a) Bentuk diferensial tarik baliknya adalah

\mavergleichskettealigndrucklinks
{\vergleichskettealigndrucklinks
{\varphi^* (  (x-y)dx-z^2dy+dz   )
}
{ =} { (u^3-u^2-v^2 ) d(u^3)  -u^{-2}v^{-2} d(u^2+v^2) +d (u^{-1}v^{-1})
}
{ =} { 3 (u^5-u^4-u^2v^2 ) du  -2u^{-1}v^{-2} du -2u^{-2}v^{-1} dv - u^{-2} v^{-1}du - u^{-1} v^{-2}dv
}
{ =} { (3 u^5-3u^4-3u^2v^2 -2u^{-1}v^{-2}- u^{-2} v^{-1} )  du + (  -2u^{-2}v^{-1} - u^{-1} v^{-2})dv
}
{ } {
}
}
{}{}{.}

b) Integral lintasannya adalah

\mavergleichskettealigndrucklinks
{\vergleichskettealigndrucklinks
{ \int_1^c (3 t^5-3t^4-3t^2t^6 -2t^{-1}t^{-6}- t^{-2} t^{-3} )dt  +(  -2t^{-2}t^{-3} - t^{-1} t^{-6}) d t^3
}
{ =} { \int_1^c   (3 t^5-3t^4-3t^8-2t^{-7} - t^{-5})dt + ( -2t^{-5} - t^{-7})3t^2 dt
}
{ =} { \int_1^c (3 t^5-3t^4-3t^8-2t^{-7} - t^{-5}   -6t^{-3} - 3 t^{-5}      ) dt
}
{ =} { \int_1^c (-3t^8 + 3 t^5-3t^4 -6t^{-3} - 4 t^{-5}  -2t^{-7}       )  dt
}
{ =} { (-{ \frac{   1 }{  3 } } t^9 + { \frac{   1 }{  2 } } t^{6}  - { \frac{   3 }{  5 } } t^5  +  3 t^{-2} +  t^{-4} + { \frac{   1 }{  3 } }  t^{-6}     ) \vert_{  1 }^{  c  }
}
}
{
\vergleichskettefortsetzungalign
{ =} { -{ \frac{   1 }{  3 } } c^9 + { \frac{   1 }{  2 } } c^{6}  - { \frac{   3 }{  5 } } c^5  +  3c^{-2} +   c^{-4} + { \frac{   1 }{  3 } }  c^{-6}
-(-{ \frac{   1 }{  3 } } + { \frac{   1 }{  2 } } - { \frac{   3 }{  5 } } +3 + 1 + { \frac{   1 }{  3 } }  )
}
{ =} { -{ \frac{   1 }{  3 } } c^9 + { \frac{   1 }{  2 } } c^{6}  - { \frac{   3 }{  5 } } c^5 + 3 c^{-2} +  c^{-4} + { \frac{   1 }{  3 } } c^{-6} - { \frac{   39 }{  10 } }  }
{ } {}
{ } {}
}{}{.}

 }



\inputaufgabeklausurloesung
{3}

{

Misalkan
\maabbdisp {\psi} { L } { M
} {}
adalah
\definitionsverweis {pemetaan diferensiabel}{}{}
antara
\definitionsverweis {manifold diferensiabel}{}{.}
Selanjutnya, misalkan
\maabbdisp {f} { M } {\R
} {}
adalah fungsi pada $M$, dan $\omega$ adalah
$k$-\definitionsverweis {bentuk}{}{}
pada $M$. Tunjukkan bahwa

\mavergleichskettedisp
{\vergleichskette
{ \psi^*(f \omega)
}
{ =} { \psi^*(f)  \psi^* \omega
}
{ } {
}
{ } {
}
{ } {
}
}
{}{}{.}

}
{

Misalkan
\mathl{P \in L}{} dan
\mathl{v_1 , \ldots , v_k \in T_PL}{.} Maka

\mavergleichskettealign
{\vergleichskettealign
{ \psi^* \left(  f \omega) (P,v_1 , \ldots , v_k  \right)
}
{ =} { (f \omega) \left(  \psi(P), T_P\psi (v_1) , \ldots , T_P\psi (v_k )       \right)
}
{ =} { f(\psi(P))   \cdot \omega \left(  \psi(P), T_P\psi (v_1) , \ldots , T_P\psi (v_k )       \right)
}
{ =} { \left(  \psi^* (f)  \right) (P)  \cdot \omega \left(  \psi(P), T_P\psi (v_1) , \ldots , T_P\psi (v_k )       \right)
}
{ =} { \left(  \psi^* (f)  \right) (P)   \cdot \left(  \psi^*( \omega)  \right)  \left(  P,v_1 , \ldots , v_k  \right)
}
}
{
\vergleichskettefortsetzungalign
{ =} { \left(  \psi^* (f)  \psi^*( \omega)  \right) \left(  P,v_1 , \ldots , v_k  \right)
}
{ } {}
{ } {}
{ } {}
}
{}{.}

 }



\inputaufgabeklausurloesung
{2}

{

Deskripsikan sebanyak mungkin
\definitionsverweis {manifold berbatas}{}{}
berdimensi satu yang
\definitionsverweis {terhubung}{}{.}

}
{  }



\inputaufgabeklausurloesung
{5}

{

Misalkan
\maabbdisp {f} { [a,b] } { \R_{+}
} {}
adalah
\definitionsverweis {fungsi yang terdiferensialkan secara kontinu}{}{.}
Kita pandang
\definitionsverweis {subgraf}{}{}
sebagai
\definitionsverweis {manifold berbatas}{}{}.
Batasnya terdiri dari grafik, interval dasar, dan kedua sisi tegak, tetapi tidak mencakup keempat titik sudut. Buktikan secara langsung berlakunya teorema Stokes untuk
\definitionsverweis {bentuk diferensial}{}{}

\mavergleichskette
{\vergleichskette
{ \omega
}
{ = }{ xdy
}
{  }{
}
{    }{
}
{   }{
}
}
{}{}{.}

}
{

Karena

\mavergleichskettedisp
{\vergleichskette
{d(xdy)
}
{ =} { dx \wedge dy
}
{ } {
}
{ } {
}
{ } {
}
}
{}{}{}
maka, di satu pihak,

\mavergleichskettedisp
{\vergleichskette
{ \int_M dx \wedge dy
}
{ =} { \int_a^b f(t) dt
}
{ =} { F(b)-F(a)
}
{ } {
}
{ } {
}
}
{}{}{}
adalah luas subgraf, dengan $F$ suatu antiturunan dari $f$. Untuk integral lintasan, batas $M$ harus diparameterkan berlawanan arah jarum jam. Interval dasarnya diparameterkan oleh
\mathl{i: t \mapsto \begin{pmatrix} t \\ 0 \end{pmatrix}}{}, dan

\mavergleichskettedisp
{\vergleichskette
{ i^*\omega
}
{ =} { i^*xdy
}
{ =} { x d0
}
{ =} { 0
}
{ } {
}
}
{}{}{,}
sehingga bagian integral lintasan ini bernilai $0$. Sisi kanan diparameterkan oleh
$i: t \mapsto \begin{pmatrix}  b  \\t   \end{pmatrix}$,
; di sini $t$ bergerak antara
\mathkor {} {0} {dan} {f(b)} {}
. Maka

\mavergleichskettedisp
{\vergleichskette
{ i^*xdy
}
{ =} { b dt
}
{ } {
}
{ } {
}
{ } {
}
}
{}{}{}
dan integral lintasan pada sisi ini adalah
\mathl{b f(b)}{.} Kontribusi sisi kiri adalah
\mathl{-af(a)}{.} Grafik diparameterkan oleh
\maabbeledisp {i} {[a,b]} {
} { t } {(t, f(t))
} {,}
. Maka

\mavergleichskettedisp
{\vergleichskette
{ i^*xdy
}
{ =} { t df
}
{ =} { tf'(t) dt
}
{ } {
}
{ } {
}
}
{}{}{.}
Dengan orientasi yang benar, integralnya adalah

\mavergleichskettealign
{\vergleichskettealign
{ - \int_a^b tf'(t)dt
}
{ =} { - ( tf(t) - F (t) ) \mid _a^b
}
{ =} { - ( bf(b)-F(b) -af(a) +F(a) )
}
{ =} { F(b)-F(a) +af(a) - bf(b)
}
{ } {
}
}
{}
{}{.}
Jadi, seluruh integral lintasan juga bernilai

\mavergleichskettealign
{\vergleichskettealign
{bf(b) - af(a)+ F(b)-F(a) +af(a) - bf(b)
}
{ =} { F(b)-F(a)
}
{ } {
}
{ } {
}
{ } {
}
}
{}
{}{.}

 }



\inputaufgabeklausurloesung
{9}

{

Buktikan teorema titik tetap Brouwer.

}
{

Untuk menyederhanakan notasi, misalkan

\mavergleichskette
{\vergleichskette
{ r
}
{ = }{ 1
}
{ }{
}
{ }{
}
{ }{
}
}
{}{}{.}  Andaikan terdapat pemetaan terdiferensialkan secara kontinu $\psi$ tanpa titik tetap. Maka selalu berlaku

\mavergleichskettedisp
{\vergleichskette
{ x
}
{ \neq} { \psi(x)
}
{ } {
}
{ } {
}
{ } {
}
}
{}{}{,}
sehingga kedua titik itu menentukan sebuah garis. Gagasannya adalah menggunakan garis ini untuk memilih
\zusatzklammer {dari keduanya} {} {}
salah satu dari kedua titik potongnya dengan sfera sebagai titik citra suatu retraksi ke batas. Dengan fungsi bantu

\mavergleichskettedisp
{\vergleichskette
{ h(x)
}
{ =} { { \frac{ \psi(x)-x }{ \Vert { \psi (x)-x} \Vert } }
}
{ } {
}
{ } {
}
{ } {
}
}
{}{}{}
kita definisikan pemetaan
\maabbdisp {\varphi} { B \left( 0 , 1 \right) } { \R^n
} {}
melalui

\mavergleichskettedisphandlinks
{\vergleichskettedisphandlinks
{ \varphi(x)
}
{ =} { x + \left(  - \left\langle  x  , h(x)  \right\rangle + \sqrt{1 + \left\langle  x  , h(x)  \right\rangle^2 - \Vert { x} \Vert^2  }  \right) \cdot h(x)
}
{ } {
}
{ } {
}
{ } {
}
}
{}{}{.}
Ekspresi di bawah tanda akar bernilai positif. Untuk

\mavergleichskette
{\vergleichskette
{ \Vert { x } \Vert
}
{ < }{ 1
}
{ }{
}
{ }{
}
{ }{
}
}
{}{}{}
hal ini jelas; sedangkan untuk

\mavergleichskette
{\vergleichskette
{ \Vert { x } \Vert
}
{ = }{ 1
}
{ }{
}
{ }{
}
{ }{
}
}
{}{}{}
terdapat titik pada sfera yang garis penghubungnya dengan titik dalam bola
\mathl{\psi(x)}{} tidak tegak lurus terhadap $x$
\zusatzklammer {ruang singgung afin pada suatu titik sfera memotong bola hanya di satu titik} {} {,}
sehingga

\mavergleichskettedisp
{\vergleichskette
{ \left\langle x , h(x) \right\rangle
}
{ \neq} { 0
}
{ } {
}
{ } {
}
{ } {
}
}
{}{}{}
berlaku. Karena fungsi akar kuadrat dan norma di luar titik nol
\definitionsverweis {terdiferensialkan secara kontinu}{}{}
bersifat demikian, maka
\mathkor {} {h(x)} {dan} {\varphi(x)} {}
merupakan pemetaan-pemetaan terdiferensialkan secara kontinu. Menurut
Soal 23.23 (Differentialgeometrie (Osnabrück 2023)),
pemetaan $\varphi$ memetakan bola ke sfera dan pembatasannya pada sfera adalah identitas. Jadi terdapat suatu
\definitionsverweis {Retraksi}{}{}
bola pejal tertutup ke batasnya, yang menurut
Teorema 23.9 (Differentialgeometrie (Osnabrück 2023))
mustahil ada.

 }



\inputaufgabeklausurloesung
{5 (3+1+1)}

{

Tinjau silinder

\mavergleichskette
{\vergleichskette
{ Z
}
{ = }{ S^1 \times \R
}
{ \subseteq }{ \R^2 \times \R
}
{    }{
}
{   }{
}
}
{}{}{}
dengan
\definitionsverweis {koneksi trivial}{}{}
pada
\definitionsverweis {bundel tangen}{}{.}
\aufzaehlungdrei{Tentukan, dalam suatu peta isometrik yang sesuai,
\definitionsverweis {persamaan diferensial geodesik}{}{}
untuk kurva geodesik $\psi$ pada $M$ yang memenuhi

\mavergleichskettedisp
{\vergleichskette
{ \psi(0)
}
{ =} { (1,0,3)
}
{ } {
}
{ } {
}
{ } {
}
}
{}{}{}
dan

\mavergleichskettedisp
{\vergleichskette
{ \psi'(0)
}
{ =} { (0,1,-4)
}
{ } {
}
{ } {
}
{ } {
}
}
{}{}{}
.
}{Selesaikan persamaan diferensial dari (1) dalam peta tersebut.
}{Carilah suatu kurva geodesik di $Z$ yang memenuhi syarat-syarat pada (1).
}

}
{

\aufzaehlungdrei{Kita tinjau peta isometrik
\zusatzklammer {inverse} {} {}
peta
\maabbeledisp {\varphi} { ]- \pi, \pi[  \times \R } { S^1 \setminus \{  (-1,0) \} \times \R
} {(t,v)} { \left(  \cos  t   , \,   \sin  t  , \,  v                 \right)
} {.}
Karena

\mavergleichskette
{\vergleichskette
{ \varphi(0,3)
}
{ = }{ (1,0,3)
}
{ }{
}
{ }{
}
{ }{
}
}
{}{}{.}
simbol-simbol Christoffelnya nol, persamaan diferensial geodesiknya adalah

\mavergleichskettedisp
{\vergleichskette
{ \gamma^{\prime \prime}
}
{ =} { 0
}
{ } {
}
{ } {
}
{ } {
}
}
{}{}{.}
Syarat-syarat yang diberikan dalam peta ini menjadi

\mavergleichskettedisp
{\vergleichskette
{ \gamma(0)
}
{ =} { (0,3)
}
{ } {
}
{ } {
}
{ } {
}
}
{}{}{}
dan

\mavergleichskettedisp
{\vergleichskette
{ \gamma'(0)
}
{ =} { (1,-4)
}
{ } {
}
{ } {
}
{ } {
}
}
{}{}{.}
}{Suatu solusi dalam peta tersebut adalah

\mavergleichskettedisp
{\vergleichskette
{ \gamma(t)
}
{ =} { (0,3) + t (1,-4)
}
{ =} { (t,3-4t)
}
{ } {
}
{ } {
}
}
{}{}{.}
}{Karena itu, kurva geodesik dengan syarat awal tersebut adalah

\mavergleichskettedisp
{\vergleichskette
{  \psi(t)
}
{ =} { \varphi( \gamma(t))
}
{ =} { \left(  \cos  t   , \,   \sin  t  , \,   3-4t                 \right)
}
{ } {
}
{ } {
}
}
{}{}{.}
}

 }
